% GENERATED FILE. Edit canonical lessons, then run npm run generate:books.
% canonical-source-hash: fnv1a64:355d1bdc2e99161d
% canonical-lessons: RU-C278-vpervyye, RU-C278-navsegda, RU-C278-nenadolgo, RU-C278-vspominat, RU-C278-zapominat, RU-R278-review-271-278

\chapter{The first time and for ever}
\label{ch:ru-the-first-time-and-for-ever}
\hlchaptermodality{278}

\begin{chapteropening}
\textbf{By the end of this chapter:} \emph{I can say it happened for the first time, for good, for a short while, and that I recall or memorize something.}

Complete the last lesson of chapter 278: I can say it happened for the first time, for good, for a short while, and that I recall or memorize something.
\end{chapteropening}

\section[\texorpdfstring{\ru{впервые} (vpervýye)}{впервые (vpervýye)}]{\ru{впервые} (vpervýye) — for the first time}

\label{lesson:RU-C278-vpervyye}



\begin{quote}
Before the new one: say the Russian for annual, then the Russian for further.
\end{quote}

\subsection*{You'll want to know: \ru{впервые}}
\textbf{\ru{впервые}} (\emph{vpervýye}) — \textquotedblleft{}for the first time\textquotedblright{}.

\begin{quote}
\ru{Я} \ru{впервые} \ru{увидел} \ru{море} \ru{в} \ru{девять} \ru{лет}.
\end{quote}

\begin{quote}
\emph{Ya vpervýye uvídel móre v dévyat let.}
\end{quote}

\begin{quote}
I saw the sea for the first time at nine.
\end{quote}

Also \emph{\ru{в} \ru{первый} \ru{раз}}.

\begin{cousinweb}
\emph{\ru{впервые}} is built on \emph{\ru{в}} and \emph{\ru{первый}} \textquotedblleft{}first\textquotedblright{}, which is inherited from Proto-Slavic \emph{pĭrvŭ} \textquotedblleft{}first\textquotedblright{} (EtymDB).
\end{cousinweb}

\begin{grammarlens}[title={What you've built}]
One more word for telling, explaining and describing at B1.
\end{grammarlens}

\subsection*{Your turn}
\begin{itemize}
\raggedright
  \item \emph{Say it:} \emph{vpervýye}
  \item \emph{Say it:} \emph{vpervýye}, once more
  \item \emph{Say it:} a sentence of your own with \emph{vpervýye}
  \item \emph{Recall:} say the Russian for annual, then the Russian for further, then say \emph{vpervýye} again
\end{itemize}

\begin{tcolorbox}[breakable,colback=teal!4,colframe=teal!35!black,title={Before you move on}]
What does \ru{впервые} mean? (\textquotedblleft{}For the first time\textquotedblright{}.) Say it once more.
\end{tcolorbox}
\section[\texorpdfstring{\ru{навсегда} (navsegdá)}{навсегда (navsegdá)}]{\ru{навсегда} (navsegdá) — for ever, for good}

\label{lesson:RU-C278-navsegda}



\begin{quote}
Before the new one: say the Russian for further, then the Russian for for the first time.
\end{quote}

\subsection*{You'll want to know: \ru{навсегда}}
\textbf{\ru{навсегда}} (\emph{navsegdá}) — \textquotedblleft{}for ever, for good\textquotedblright{}.

\begin{quote}
\ru{Он} \ru{уехал} \ru{из} \ru{города} \ru{навсегда}.
\end{quote}

\begin{quote}
\emph{On uyékhal iz góroda navsegdá.}
\end{quote}

\begin{quote}
He left the town for good.
\end{quote}

\emph{\ru{Раз} \ru{и} \ru{навсегда}} — once and for all.

\begin{cousinweb}
\emph{\ru{навсегда}} is built on \emph{\ru{на}} and \emph{\ru{всегда}} \textquotedblleft{}always\textquotedblright{}, itself built on \emph{\ru{весь}} \textquotedblleft{}all\textquotedblright{}; EtymDB lists no parent for \emph{\ru{всегда}}.
\end{cousinweb}

\begin{grammarlens}[title={What you've built}]
One more word for telling, explaining and describing at B1.
\end{grammarlens}

\subsection*{Your turn}
\begin{itemize}
\raggedright
  \item \emph{Say it:} \emph{navsegdá}
  \item \emph{Say it:} \emph{navsegdá}, once more
  \item \emph{Say it:} a sentence of your own with \emph{navsegdá}
  \item \emph{Recall:} say the Russian for further, then the Russian for for the first time, then say \emph{navsegdá} again
\end{itemize}

\begin{tcolorbox}[breakable,colback=teal!4,colframe=teal!35!black,title={Before you move on}]
What does \ru{навсегда} mean? (\textquotedblleft{}For ever, for good\textquotedblright{}.) Say it once more.
\end{tcolorbox}
\section[\texorpdfstring{\ru{ненадолго} (nenadólgo)}{ненадолго (nenadólgo)}]{\ru{ненадолго} (nenadólgo) — for a short while, not for long}

\label{lesson:RU-C278-nenadolgo}



\begin{quote}
Before the new one: say the Russian for for the first time, then the Russian for for ever.
\end{quote}

\subsection*{You'll want to know: \ru{ненадолго}}
\textbf{\ru{ненадолго}} (\emph{nenadólgo}) — \textquotedblleft{}for a short while, not for long\textquotedblright{}.

\begin{quote}
\ru{Она} \ru{зашла} \ru{ненадолго} \ru{и} \ru{скоро} \ru{ушла}.
\end{quote}

\begin{quote}
\emph{Oná zashlá nenadólgo i skóro ushlá.}
\end{quote}

\begin{quote}
She dropped in for a short while and soon left.
\end{quote}

The opposite is \emph{\ru{надолго}} — for a long time.

\begin{cousinweb}
\emph{\ru{ненадолго}} is built on \emph{\ru{не}}, \emph{\ru{на}} and \emph{\ru{долго}} \textquotedblleft{}for long\textquotedblright{}, the adverb of \emph{\ru{долгий}} \textquotedblleft{}long\textquotedblright{}, which is inherited from Proto-Slavic \emph{dĭlgŭ} (EtymDB).
\end{cousinweb}

\begin{grammarlens}[title={What you've built}]
One more word for telling, explaining and describing at B1.
\end{grammarlens}

\subsection*{Your turn}
\begin{itemize}
\raggedright
  \item \emph{Say it:} \emph{nenadólgo}
  \item \emph{Say it:} \emph{nenadólgo}, once more
  \item \emph{Say it:} a sentence of your own with \emph{nenadólgo}
  \item \emph{Recall:} say the Russian for for the first time, then the Russian for for ever, then say \emph{nenadólgo} again
\end{itemize}

\begin{tcolorbox}[breakable,colback=teal!4,colframe=teal!35!black,title={Before you move on}]
What does \ru{ненадолго} mean? (\textquotedblleft{}For a short while, not for long\textquotedblright{}.) Say it once more.
\end{tcolorbox}
\section[\texorpdfstring{\ru{вспоминать} (vspominát)}{вспоминать (vspominát)}]{\ru{вспоминать} (vspominát) — to recall, to remember}

\label{lesson:RU-C278-vspominat}



\begin{quote}
Before the new one: say the Russian for for ever, then the Russian for for a short while.
\end{quote}

\subsection*{You'll want to know: \ru{вспоминать}}
\textbf{\ru{вспоминать}} (\emph{vspominát}) — \textquotedblleft{}to recall, to remember\textquotedblright{}.

\begin{quote}
\ru{Бабушка} \ru{часто} \ru{вспоминает} \ru{свою} \ru{деревню}.
\end{quote}

\begin{quote}
\emph{Bábushka chásto vspomináyet svoyú derévnyu.}
\end{quote}

\begin{quote}
Grandma often recalls her village.
\end{quote}

Imperfective; perfective \emph{\ru{вспомнить}} (\emph{vspómnit}). It takes the accusative or \emph{\ru{о}} + prepositional. \emph{\ru{Помнить}} is to keep in memory; \emph{\ru{вспоминать}} is to bring something back to mind.

\begin{cousinweb}
\emph{\ru{вспоминать}} is built on \emph{\ru{вс}-} \textquotedblleft{}up\textquotedblright{} and \emph{\ru{помнить}} \textquotedblleft{}to remember\textquotedblright{}, from the Slavic root of \emph{\ru{память}} \textquotedblleft{}memory\textquotedblright{}; EtymDB lists no parent for the compound.
\end{cousinweb}

\begin{grammarlens}[title={What you've built}]
One more word for telling, explaining and describing at B1.
\end{grammarlens}

\subsection*{Your turn}
\begin{itemize}
\raggedright
  \item \emph{Say it:} \emph{vspominát}
  \item \emph{Say it:} \emph{vspominát}, once more
  \item \emph{Say it:} a sentence of your own with \emph{vspominát}
  \item \emph{Recall:} say the Russian for for ever, then the Russian for for a short while, then say \emph{vspominát} again
\end{itemize}

\begin{tcolorbox}[breakable,colback=teal!4,colframe=teal!35!black,title={Before you move on}]
What does \ru{вспоминать} mean? (\textquotedblleft{}To recall, to remember\textquotedblright{}.) Say it once more.
\end{tcolorbox}
\section[\texorpdfstring{\ru{запоминать} (zapominát)}{запоминать (zapominát)}]{\ru{запоминать} (zapominát) — to memorize, to commit to memory}

\label{lesson:RU-C278-zapominat}



\begin{quote}
Before the new one: say the Russian for for a short while, then the Russian for to recall.
\end{quote}

\subsection*{You'll want to know: \ru{запоминать}}
\textbf{\ru{запоминать}} (\emph{zapominát}) — \textquotedblleft{}to memorize, to commit to memory\textquotedblright{}.

\begin{quote}
\ru{Я} \ru{легко} \ru{запоминаю} \ru{имена}.
\end{quote}

\begin{quote}
\emph{Ya legkó zapomináyu imená.}
\end{quote}

\begin{quote}
I remember names easily.
\end{quote}

Imperfective; perfective \emph{\ru{запомнить}} (\emph{zapómnit}): \emph{\ru{я} \ru{запомнил} \ru{дорогу}} — I memorized the way.

\begin{cousinweb}
\emph{\ru{запоминать}} is built on \emph{\ru{за}-} and \emph{\ru{помнить}} \textquotedblleft{}to remember\textquotedblright{}, from the Slavic root of \emph{\ru{память}} \textquotedblleft{}memory\textquotedblright{}: to fix in memory. EtymDB lists no parent for the compound.
\end{cousinweb}

\begin{grammarlens}[title={What you've built}]
One more word for telling, explaining and describing at B1.
\end{grammarlens}

\subsection*{Your turn}
\begin{itemize}
\raggedright
  \item \emph{Say it:} \emph{zapominát}
  \item \emph{Say it:} \emph{zapominát}, once more
  \item \emph{Say it:} a sentence of your own with \emph{zapominát}
  \item \emph{Recall:} say the Russian for for a short while, then the Russian for to recall, then say \emph{zapominát} again
\end{itemize}

\begin{tcolorbox}[breakable,colback=teal!4,colframe=teal!35!black,title={Before you move on}]
What does \ru{запоминать} mean? (\textquotedblleft{}To memorize, to commit to memory\textquotedblright{}.) Say it once more.
\end{tcolorbox}
\section[(dialogue)]{Telling what happened — a review}

\label{lesson:RU-R278-review-271-278}



\begin{quote}
Say the Russian for suddenly and instantly.
\end{quote}

\subsection*{Your turn}
Say these aloud:

\begin{itemize}
\raggedright
  \item suddenly, instantly, at once, immediately, an instant
  \item in advance, the day before, before, soon after, later
  \item to happen, to go on, to last, to stop, to come
  \item gradually, at the same time, at first, later on, at the very end
  \item to drag on, to put off, to be delayed, to happen again, to coincide
  \item a decade, a period, a stage, an interval, a deadline
  \item previous, following, next, annual, further
  \item for the first time, for ever, for a short while, to recall, to memorize
\end{itemize}

\begin{tcolorbox}[breakable,colback=teal!4,colframe=teal!35!black,title={Before you move on}]
Suddenly? (\textbf{\ru{внезапно}}.) Instantly? (\textbf{\ru{мгновенно}}.) At once? (\textbf{\ru{тотчас}}.) Immediately? (\textbf{\ru{немедленно}}.) An instant? (\textbf{\ru{мгновение}}.) In advance? (\textbf{\ru{заранее}}.) The day before? (\textbf{\ru{накануне}}.) Before? (\textbf{\ru{прежде}}.) Soon after? (\textbf{\ru{вскоре}}.) Later? (\textbf{\ru{спустя}}.) To happen? (\textbf{\ru{происходить}}.) To go on? (\textbf{\ru{продолжаться}}.) To last? (\textbf{\ru{длиться}}.) To stop? (\textbf{\ru{прекращаться}}.) To come? (\textbf{\ru{наступать}}.) Gradually? (\textbf{\ru{постепенно}}.) At the same time? (\textbf{\ru{одновременно}}.) At first? (\textbf{\ru{поначалу}}.) Later on? (\textbf{\ru{впоследствии}}.) At the very end? (\textbf{\ru{напоследок}}.) To drag on? (\textbf{\ru{затягиваться}}.) To put off? (\textbf{\ru{откладывать}}.) To be delayed? (\textbf{\ru{задерживаться}}.) To happen again? (\textbf{\ru{повторяться}}.) To coincide? (\textbf{\ru{совпадать}}.) A decade? (\textbf{\ru{десятилетие}}.) A period? (\textbf{\ru{период}}.) A stage? (\textbf{\ru{этап}}.) An interval? (\textbf{\ru{промежуток}}.) A deadline? (\textbf{\ru{срок}}.) Previous? (\textbf{\ru{предыдущий}}.) Following? (\textbf{\ru{последующий}}.) Next? (\textbf{\ru{очередной}}.) Annual? (\textbf{\ru{ежегодный}}.) Further? (\textbf{\ru{дальнейший}}.) For the first time? (\textbf{\ru{впервые}}.) For ever? (\textbf{\ru{навсегда}}.) For a short while? (\textbf{\ru{ненадолго}}.) To recall? (\textbf{\ru{вспоминать}}.) To memorize? (\textbf{\ru{запоминать}}.)
\end{tcolorbox}
